% Generated draft: requires mathematical and source review.
\documentclass[11pt]{article}
\usepackage[margin=25mm]{geometry}
\usepackage{fontspec}
\setmainfont{Latin Modern Roman}
\usepackage{amsmath,amssymb,mathtools,unicode-math}
\setmathfont{Latin Modern Math}
\usepackage{enumitem,needspace,xurl,hyperref}
\setlength{\parindent}{0pt}
\setlength{\parskip}{5pt}
\setlength{\emergencystretch}{3em}
\newcommand{\N}{\mathbb{N}}
\newcommand{\Z}{\mathbb{Z}}
\newcommand{\Q}{\mathbb{Q}}
\newcommand{\R}{\mathbb{R}}
\newcommand{\C}{\mathbb{C}}
\newcommand{\F}{\mathbb{F}}
\newcommand{\A}{\mathbb{A}}
\newcommand{\PP}{\mathbb{P}}
\newcommand{\E}{\mathbb{E}}
\newcommand{\eps}{\varepsilon}
\newcommand{\dd}{\,\mathrm{d}}
\begin{document}
\section*{UnsolvedMath ID: 30003974}
\textbf{Source problem code:} catalogue entry
\subsection{Definitions and mathematical statement}
For a standard single-point interpretation, fix an algebraically closed field \(k\). Let \(X\) be a smooth projective integral variety over \(k\) with \(\dim X\geq 1\), let \(L\) be an ample line bundle on \(X\), and let \(x\in X(k)\). The Seshadri constant of \(L\) at \(x\) is
\[
\varepsilon(L;x)=\inf_{\substack{C\subseteq X\text{ an integral curve}\\x\in C}}\frac{L\cdot C}{\operatorname{mult}_x(C)},
\]
where \(L\cdot C\) is the degree of the pullback of \(L\) to the normalization of \(C\), and \(\operatorname{mult}_x(C)\) is the multiplicity of \(C\) at \(x\).

Under this interpretation, the question is whether there exist \(X,L,x\) satisfying these hypotheses such that
\[
\varepsilon(L;x)\in\R\setminus\Q.
\]
This is a conditional formulation: the supplied record does not specify the ground field, its characteristic, the dimension or admissible class of varieties, or whether single-point Seshadri constants are intended. The assumptions above describe a standard setting but are not explicitly supplied hypotheses. The record also does not distinguish single-point constants from multipoint or global variants, or specify the permitted positivity and coefficient class. Consequently, the precise intended scope cannot be uniquely recovered.
\subsection{Short English statement}
Can a Seshadri constant be irrational? In the standard single-point setting, this asks whether the infimum of the ratios of an ample line bundle's degree on a curve to that curve's multiplicity at a fixed point can be irrational. The source leaves the field, geometric setting, and variant of the constant unspecified.
\subsection{Sources}
\begin{itemize}
\item Supplied local UnsolvedMath registry record, permanent ID \(30003974\), title: \emph{Irrationality of Seshadri Constants}; collection: Oberwolfach Reports Open Problems; exact source problem code: \texttt{catalogue entry}. The supplied original statement is ``Can Seshadri constants be irrational?'' No individual proposer is identified.
\item Supplied source citation and literature source labeled ``OWR report'': \emph{Mini-Workshop: Asymptotic Invariants of Homogeneous Ideals (2019)}, Oberwolfach Reports, DOI 10.4171/owr/2018/45, pp.\ 2739--2767. \url{https://doi.org/10.4171/owr/2018/45}. The record gives proposed year 2018; its citation uses 2019, while its DOI and year tag refer to 2018. This dating discrepancy is unresolved.
\item Supplied literature source: \emph{Seshadri constants in characteristic zero}. \url{https://www.sciencedirect.com/science/article/abs/pii/S0007449722001130}.
\item Supplied literature source: \emph{Seshadri functions on abelian surfaces}. \url{https://doi.org/10.1016/j.aim.2022.108827}.
\item Literature assessment supplied with exact check date \texttt{2026-08-21}: the record labels the question open and reports that later work computes broad classes, including abelian surfaces, without producing an unconditional irrational example. This is supplied dated curation derived from the OWR curation, not an independently verified literature assessment. Its time-sensitive status and cited claims require independent verification.
\end{itemize}
\end{document}
